\section{Construction Quality}
\label{sec:construction_quality}
\begin{table}[t]
\centering
\small
\caption{\textbf{Final dataset statistics.} Number of source questions and generated questions by transformation.}
\label{tab:stats}
\begin{tabular}{@{}lrrrr@{}}
\toprule
 & & \multicolumn{2}{c}{\textbf{KMMLU}} & \\
\cmidrule(lr){3-4}
 & \textbf{CLIcK} & \textbf{Pro} & \textbf{Redux} & \textbf{Total} \\
\midrule
Source questions    & 66  & 53 & 46 & 165 \\
\midrule
A1$\rightarrow$A2   & 29  & 20 & 29 & 78  \\
A1$\rightarrow$A3   & 24  & 15 & 12 & 51  \\
A2$\rightarrow$A3   & 77  & 55 & 26 & 158 \\
\midrule
Generated questions & 130 & 90 & 67 & 287 \\
\bottomrule
\end{tabular}
\end{table}
\begin{table}[t]
\centering
\small
\caption{\textbf{Stage-wise yield.}
Indented rates are relative to the preceding unindented count; the overall automated-stage pass rate is relative to all generated candidates.}
\label{tab:yield}
\begin{tabular}{@{}lrr@{}}
\toprule
\textbf{Stage} & \textbf{Count} & \textbf{Rate} \\
\midrule
\multicolumn{3}{@{}l}{\textit{Automated stage}} \\
\cmidrule{1-3}
Candidates generated                    & 1,964 & --     \\
\quad All eleven criteria \textsc{pass} & 627   & 31.9\% \\
\quad Proxy answers incorrectly         & 760   & 38.7\% \\
\quad Both (pass before revision)       & 294   & 15.0\% \\
\cmidrule{1-3}
Candidates revised                      & 1,662 & --     \\
\quad Pass after self-revision          & 329   & 19.8\% \\
\cmidrule{1-3}
Pass the automated stage                & 623   & 31.7\% \\
\midrule[\heavyrulewidth]
\multicolumn{3}{@{}l}{\textit{Human validation}} \\
\cmidrule{1-3}
Candidates validated                    & 301   & --     \\
\quad Pass directly                     & 252   & 83.7\% \\
\cmidrule{1-3}
Candidates revised                      & 49    & --     \\
\quad Pass after human-guided revision  & 35    & 71.4\% \\
\cmidrule{1-3}
Final questions                         & 287   & --     \\
\bottomrule
\end{tabular}
\vspace{-0.1in}
\end{table}
% \input{Tables/Table_criteria}
Next, we examine whether \name{} produces valid knowledge-matched contrasts and how much each validation and revision stage contributes.

\subsection{Experimental Setup}
\label{sec:construction_setup}

We use 243 seeds, 78 at A1 and 165 at A2, which form 295 seed--target pairs: 78 A1$\rightarrow$A2, 52 A1$\rightarrow$A3, and 165 A2$\rightarrow$A3.
We use GPT-5.5 to construct anchor bundles, generate and revise candidates, and audit them, and GPT-5-nano as the proxy model.
For each seed--target pair, we first generate three candidates and, if fewer than two pass auditing and screening, generate six more in one additional round.
We additionally examine how seed difficulty affects the results in Appendix~\ref{app:easy}.

Human validation proceeds in two batches.
In the first batch (164 pairs), up to two passing candidates per pair are validated.
Two annotators judge all five criteria, and two additional annotators double-check H4 and H5.
In the second batch (131 pairs), 111 pairs have a passing candidate; because of annotation budget and time constraints, we validate one candidate for 70 of them, each judged by one of three annotators. 
All annotators are native Korean speakers who work as professional data annotators and follow written guidelines calibrated on a pilot set.
Table~\ref{tab:stats} summarizes the final dataset of 287 questions built from 165 source questions.

\subsection{Construction Validity and Refinement}
\label{sec:construction_results}

\paragraph{Yield and validity.}
We first examine how many candidates survive each stage and whether the surviving ones are valid contrasts (Table~\ref{tab:yield}).
\name{} generates 1,964 candidates for the 295 pairs, of which 623 (31.7\%) pass the automated stage.
At least one candidate passes for 243 pairs (82.4\%), and human validation covers 202 of them with 301 candidates.
Of these candidates, 252 (83.7\%) pass human validation directly, and 35 more pass after human-guided revision, giving 287 final questions.
The two criteria that define a knowledge-matched contrast hold for nearly all validated candidates: knowledge sufficiency (H4) for 98.0\% and target-level structure (H5) for 95.7\%.
Although only a minority of candidates survive the automated stage, those that do largely satisfy the intended controls.

\begin{table*}[t]
\centering
\small
\caption{\textbf{Closed-book accuracy (\%) by application level.}
\textit{Matched} covers the 30 seeds that have both an A2 and an A3 contrast (40 and 43 questions), averaged over seeds.
The numbers of A2 and A3 questions are 29 and 101 for CLIcK, 20 and 70 for KMMLU-Pro, and 29 and 38 for KMMLU-Redux.}
\label{tab:matched}
\begin{tabular}{@{}l|rrrrrrrr|rr@{}}
\toprule
 & & & \multicolumn{4}{c}{\textbf{KMMLU}} & & & & \\
\cmidrule(lr){4-7}
 & \multicolumn{2}{c}{\textbf{CLIcK}} & \multicolumn{2}{c}{\textbf{Pro}} & \multicolumn{2}{c}{\textbf{Redux}} & \multicolumn{2}{c}{\textbf{All}} & \multicolumn{2}{c}{\textbf{Matched}} \\
\cmidrule(lr){2-3} \cmidrule(lr){4-5} \cmidrule(lr){6-7} \cmidrule(lr){8-9} \cmidrule(l){10-11}
\textbf{Model} & \textbf{A2} & \textbf{A3} & \textbf{A2} & \textbf{A3} & \textbf{A2} & \textbf{A3} & \textbf{A2} & \textbf{A3} & \textbf{A2} & \textbf{A3} \\
\midrule
K-EXAONE          & 65.5  & 46.5 & 50.0  & 37.1 & 34.5 & 57.9 & 50.0 & 45.5 & 46.7 & 40.0 \\
A.X K1            & 72.4  & 40.6 & 45.0  & 27.1 & 48.3 & 47.4 & 56.4 & 37.3 & 51.7 & 21.7 \\
Solar Open        & 55.2  & 29.7 & 30.0  & 21.4 & 34.5 & 36.8 & 41.0 & 28.2 & 38.3 & 23.3 \\
\midrule
Aya Expanse 32B   & 31.0  & 26.7 & 25.0  & 21.4 & 20.7 & 18.4 & 25.6 & 23.4 & 33.3 & 20.0 \\
\midrule
GLM-4.7           & 75.9  & 67.3 & 55.0  & 48.6 & 58.6 & 63.2 & 64.1 & 60.3 & 65.0 & 56.7 \\
Qwen3.5           & 82.8  & 72.3 & 75.0  & 68.6 & 72.4 & 60.5 & 76.9 & 68.9 & 76.7 & 58.3 \\
\midrule
Claude Sonnet 5.5 & 100.0 & 86.1 & 100.0 & 77.1 & 79.3 & 78.9 & 92.3 & 81.8 & 93.3 & 81.7 \\
GPT-6.1 Sol       & 96.6  & 94.1 & 100.0 & 92.9 & 79.3 & 86.8 & 91.0 & 92.3 & 93.3 & 93.3 \\
\bottomrule
\end{tabular}
\vspace{-0.1in}
\end{table*}

\paragraph{Effect of the anchor bundle.}
We examine what the anchor bundle contributes by comparing candidates generated with and without it for the same pairs.
In both cases, the generator is instructed to rely only on the knowledge of the seed, but without the bundle, 20.4\% of candidates require domain knowledge beyond it (L4), compared with 1.9\% with the bundle.
These candidates introduce knowledge that the seed never required, such as a rule that the bundle's scope note explicitly excludes.
An instruction to preserve knowledge is thus not enough, and the knowledge must be specified explicitly.

\paragraph{Automated defects and self-revision.}
Turning to the validation stages, we examine where candidates fail in the automated stage and how much self-revision helps (Table~\ref{tab:criteria}).
Before revision, only 15.0\% of candidates pass the automated stage.
L1--L5 each fail for at most 2.2\% of candidates, whereas L6--L11 each fail for 6.6--33.5\%. 
Self-revision more than doubles the number of passing candidates, from 294 to 623, although only 19.8\% of the 1,662 revised candidates pass.
It resolves 64--100\% of the failures on each criterion but also introduces new ones, most often on shortcut resistance (L8) and anchor leakage (L9).
The main bottleneck is thus not basic validity but imposing the intended knowledge and application demands.

\paragraph{Residual defects and human-guided revision.}
Finally, we examine the defects that remain after the automated stage.
Although all 301 validated candidates satisfy every LLM criterion, annotators find at least one defect in 49 (16.3\%).
These defects are mostly isolated, with 37 of the 49 failing a single criterion, and readability (H1) is the most frequent.
Human-guided revision recovers 35 of the 49 candidates (71.4\%) and resolves 77--100\% of the failures on each criterion, including all six on knowledge sufficiency (H4).
Of the 14 candidates that are discarded, 8 fail on a criterion that the revision newly introduced.
Automated auditing alone is therefore not sufficient, and revised questions need to be validated again.